This paper was enlarged to become the geometric foundation of the series.  The
division of labour between Paper~0 and Paper~I is fixed by the authoritative
amendment `governance/00b-paper-0-geometric-foundation-amendment.md`, and this
section records the resulting interface without restating either paper's
theorems.

\subsection{What Paper 0 establishes}

The proved spine of this paper is
\[
\begin{gathered}
  \text{left and right comb grammars}
  \longrightarrow
  \text{slot-$(1)$/slot-$(2)$ one-hole maps},\\
  \text{distributive expansion}
  \longrightarrow
  \text{same operator, different charges},\\
  \text{arithmetic motions on an AES}
  \longrightarrow
  \text{point path in } \mathfrak E_0,\\
  \text{projective operator}
  \longrightarrow
  \text{pole, fixed points, ripple pencil},\\
  \text{affine cocycles}
  \longrightarrow
  \text{relative defect}
  \longrightarrow
  \text{ACS weighted area},\\
  \text{contact form}
  \longrightarrow
  \text{non-integrable horizontal distribution}
  \longrightarrow
  \text{curvature } \mu\lambda,\\
  \Aff(1,K)=\Stab(\infty)
  \subset
  \PGL_2(K).
\end{gathered}
\]
The paper also supplies four laboratories: the basic hyperbolic model
$\mathfrak E_0$, the one-puncture model $\mathfrak E_1$, the pulled-back
horocycle pencil, and the three-dimensional contact model.  The finite
geometric-series and continued-fraction examples are exact at every truncation;
their infinite values are asserted only under the stated formal or analytic
convergence hypotheses.

Two statements in the paper are not theorems about AES and are labelled as
proposals where they occur: the descriptive name *ripple geometry* for the dual
reading of Section~\ref{sec:p0-ripples}, and the language of *tearing* of
Section~\ref{sec:p0-acs-tearing} together with the punctured reading of
Section~\ref{sec:p0-holed}.  Both rest on proved computations --- the
line-to-circle theorem, the weighted torsion--Stokes theorem, and the contact
curvature --- and neither claims a new isomorphism class of geometry.

\subsection{What remains in Paper I}

Paper~I begins one level earlier and keeps the history layer.  It retains
responsibility for:

\begin{enumerate}
  \item the intrinsic dependency-poset classification of all sequential binary
  trees, including mixed slot words;

  \item marked planar sequential trees and bounded/free marked spinal
  histories, with operand-slot chirality, mirror, temporal reversal, and path
  inverse kept distinct;

  \item projective evaluation of chronological marked histories, including
  ordinary-domain data and the affine sector;
  
  \item the $q=4$ Hecke sublanguage and its history-level relations;

  \item the regular-zero theorems, the complete splitting theorem, the
  foundational definition of a singular arithmetic expression space, the
  verified isolated-zero model, and the smooth parameter-family zero-set lemma;

  \item the interfaces to Papers~II, III, and~IV.
\end{enumerate}

Paper~I states the definition of a regular AES and its canonical frame as a
concise imported interface, because its zero theorems are stated in those terms;
everything else that moved here is referenced rather than repeated.  The earlier
version of that interface, which ran in the opposite direction, is superseded.

\subsection{The dependency architecture}

The series dependency is
\[
  \boxed{
  \text{Paper 0}
  \longrightarrow
  \text{Paper I}
  \longrightarrow
  \{\text{Papers II, III, IV}\}.
  }
\]
The first arrow exports the following interfaces:

\begin{center}
\renewcommand{\arraystretch}{1.3}
\begin{tabular}{ll}
\toprule
Paper 0 object & use in later papers \\
\midrule
comb and slot conventions & operand-slot chirality on sequential spines \\
distributive expansion & comparison of expansions with equal operator \\
regular AES and canonical frame & imported definition for zero theory \\
arithmetic motions and $\mathfrak E_0$ & model for analysis and for the flow \\
one-puncture model $\mathfrak E_1$ & first punctured instance \\
point path & bounded marked-history evaluation in the affine sector \\
ripple pencil and pole & motivation for bilateral projective continuation \\
affine cocycles & discrete order defect \\
ACS and relative torsion & global torsion and its measures \\
contact form and horizontal curvature & horizontal calculus and function theory \\
tearing and punctured spaces & vocabulary for singular and resource questions \\
\bottomrule
\end{tabular}
\end{center}

Papers~II and~III import the regular-AES definition, the canonical frame, the
basic hyperbolic model, the contact model, and the elementary bilateral
operators from this paper; they do not redefine them.  Paper~IV imports the
history layer from Paper~I and the affine torsion from this paper.

\subsection{Open extensions}

Several genuine problems remain, and they are recorded here rather than
advertised as results.

\begin{enumerate}
  \item Construct a projective analogue of the affine arithmetic expression
  space in which the pole and the ripple pencil are intrinsic rather than chosen
  through one upper-half-plane chart, and decide which path or ripple data
  descend through relations among history words.

  \item Decide whether the contact structure of Section~\ref{sec:p0-contact}
  admits a canonical compatible structure beyond the horizontal distribution,
  and if so which.  Neither a complex structure nor a horizontal metric is
  determined by the contact form alone.

  \item Determine whether the choice of punctures in Section~\ref{sec:p0-holed}
  can be made canonical, and whether a punctured model carries a ripple
  structure that survives at the puncture; the three registered programmes of
  Subsection~\ref{subsec:p0-holed-open} are the concrete routes.

  \item Extend the one-hole analysis to multi-wire expression networks while
  retaining ordinary-domain, pole, and fixed-point information, and develop a
  controlled theory of infinite arithmetic expressions that separates formal,
  analytic, projective, and computational convergence.

  \item Determine whether the non-linear algebraic structures proposed in
  Section~\ref{sec:p0-beyond-linear} can be defined intrinsically for AES, so
  that the two proved obstructions of that section are replaced by a positive
  description rather than a refusal.
\end{enumerate}

None of these problems is solved merely by the matrix representation or by the
existence of the contact form.  The stable conclusion of the present paper is
more elementary: the four arithmetic operations act as motions of one space,
their order defect is measured in the commutativization of that space, and the
contact structure of the same defect is what prevents the two motions from being
calibrated independently.
